\chapter{Абетка Морзе: що ми знаємо про мову, якою говорять \nt{GLUT}- і \nt{GABA}-нейрони}
\chaptermark{Абетка Морзе}
\label{sec:code}

\section{Алфавіт з одного знака}

В абетці Морзе два знаки, крапка й тире, і паузи трьох довжин між ними. У нейрона, як ми бачили в розділі~\ref{sec:neurontypes}, алфавіт ще бідніший: знак один. Імпульс триває близько мілісекунди, і всі імпульси одного нейрона однакові за висотою й формою. Тому все повідомлення --- в паузах, тобто в послідовності моментів $t_1,t_2,t_3,\dots$ Це абетка без тире, у якій зміст несе лише ритм крапок (рис.~\ref{fig:morse}).

\begin{figure}[t]
\centering
\begin{tikzpicture}[>=Stealth,x=0.08cm,y=1cm]
% Morse letter: dot, dash, dot
\node[left,font=\small] at (-6,2.55) {Морзе:};
\fill[black] (0,2.45) rectangle (6,2.65);
\fill[black] (14,2.45) rectangle (32,2.65);
\fill[black] (40,2.45) rectangle (46,2.65);
\node[right,font=\small,gray!40!black] at (52,2.55) {крапка, тире, крапка --- зміст у тривалостях і паузах};
% stimulus onset
\node[left,font=\small] at (-6,0.4) {нейрон:};
\draw[->,thick,green!45!black] (0,-0.55) -- (0,-0.05);
\node[below,font=\scriptsize,green!45!black] at (0,-0.55) {подразник};
\draw[gray!70!black] (0,0) -- (152,0) node[right,black] {\small $t$};
% spikes: single ones blue, the burst red
\foreach \s in {14,26,41,88,112,131}{\draw[very thick,blue!60!black] (\s,0) -- (\s,0.8);}
\foreach \s in {60,63,66,69}{\draw[very thick,red!70!black] (\s,0) -- (\s,0.8);}
% latency
\draw[<->,thick] (0,1.1) -- node[above,font=\scriptsize] {$t_1$: час першого імпульсу} (14,1.1);
% burst
\draw[decorate,decoration={brace,amplitude=4pt},red!70!black] (58,0.95) -- (71,0.95) node[midway,above=4pt,font=\scriptsize] {пачка --- «тире»};
% counting windows
\foreach \a/\b/\n in {0/50/3,50/100/5,100/150/2}{
  \draw[decorate,decoration={brace,amplitude=4pt,mirror},gray!50!black] (\a+1,-0.2) -- (\b-1,-0.2) node[midway,below=4pt,font=\small] {$n=\n$};}
\node[font=\scriptsize,gray!40!black] at (75,-1.05) {частотний код: кількість імпульсів $n$ у кожному вікні};
\foreach \x in {0,50,100,150}{\draw[gray!60!black] (\x,-0.06) -- (\x,0.06);}
\end{tikzpicture}
\caption{Абетка Морзе й абетка нейрона. Ту саму послідовність імпульсів можна прочитати по-різному: порахувати імпульси у вікнах по $50\ms$ (підрозділ~\ref{sec:rate}), виміряти затримку $t_1$~\eqref{eq:latency} або помітити пачку --- «тире»~\eqref{eq:burst}.}
\label{fig:morse}
\end{figure}

Які паузи взагалі можливі? Знизу їх обмежує несприйнятливість: після імпульсу нейрон близько мілісекунди не може спрацювати знову, тож більше кількох сотень імпульсів за секунду не буває. Згори не обмежує ніщо: нейрон може мовчати хвилинами. У \nt{GLUT}-нейронів кори частоти лежать у межах від часток імпульсу до кількох десятків імпульсів за секунду, і більшість нейронів більшу частину часу працює біля нижньої межі.

У розділах~\ref{sec:glutcomputer} і~\ref{sec:gaba} ми мимохідь приймали то один, то інший спосіб запису чисел імпульсами. Тепер запитаймо прямо: якою мовою \nt{GLUT}- і \nt{GABA}-нейрони говорять одне з одним? Повної відповіді немає, цю мову не розшифровано. Але дещо про неї відомо твердо, і майже все відоме можна отримати, міркуючи як інженер зв'язку: пропускна здатність каналу, шум, приймач, ціна передачі.

\section{Скільки можна передати}

Почнімо з верхньої межі. Нехай отримувач розрізняє моменти імпульсів із точністю $\Delta t$: зсуви, менші за $\Delta t$, для нього нерозрізненні. Розріжмо час на клітинки довжини $\Delta t$, менші за час несприйнятливості, тож у кожній клітинці або один імпульс, або жодного (рис.~\ref{fig:cells}). За середньої частоти $r$ імпульс потрапляє в клітинку з імовірністю $p=r\Delta t$. Найбільше інформації потік несе, коли клітинки незалежні, і тоді кожна клітинка дає ентропію $-p\log_2p-(1-p)\log_2(1-p)\approx p\log_2(e/p)$ біт при $p\ll1$. Поділивши на $\Delta t$, отримаємо граничну швидкість передачі та її частку на один імпульс~\cite{MacKay1952}:
\begin{empheq}[box=\keybox]{equation}
R_{\max} \;\approx\; r\,\log_2\frac{e}{r\,\Delta t}\ \ \text{біт/с},
\qquad
\frac{R_{\max}}{r} \;\approx\; \log_2\frac{e}{r\,\Delta t}\ \ \text{біт на імпульс}.
\label{eq:spikeentropy}
\end{empheq}
При $r=10$ імпульсів за секунду й $\Delta t=1\ms$ це близько восьми біт на імпульс і вісімдесяти біт за секунду.

\begin{figure}[t]
\centering
\begin{tikzpicture}[>=Stealth,x=0.5cm,y=1cm]
% time cut into cells of width Delta t; a cell holds one impulse or none
\foreach \k in {0,...,23}{\draw[gray!55] (\k,0) rectangle (\k+1,0.7);}
\foreach \k in {2,7,8,15,21}{
  \fill[red!10] (\k,0) rectangle (\k+1,0.7);
  \draw[very thick,red!70!black] (\k+0.5,0.05) -- (\k+0.5,0.65);}
\foreach \k in {0,...,23}{
  \pgfmathtruncatemacro{\bit}{(\k==2)||(\k==7)||(\k==8)||(\k==15)||(\k==21)}
  \ifnum\bit=1 \def\bitcol{red!70!black}\else \def\bitcol{gray!60!black}\fi
  \node[font=\scriptsize,text=\bitcol] at (\k+0.5,-0.25) {\bit};}
\draw[<->,gray!60!black] (0,0.95) -- (1,0.95) node[midway,above,font=\scriptsize,black] {$\Delta t$};
\node[font=\small,anchor=east] at (-0.3,0.35) {імпульси};
\node[font=\small,anchor=east] at (-0.3,-0.25) {біти};
\draw[decorate,decoration={brace,amplitude=5pt,mirror},gray!60!black] (0,-0.55) -- (24,-0.55)
  node[midway,below=6pt,font=\scriptsize,black,align=center] {отримувач, який лише рахує: «$n=5$»;\\ отримувач, який бачить клітинки: які саме $5$ із $24$ --- це $\log_2\binom{24}{5}\approx15$ біт};
\end{tikzpicture}
\caption{Звідки береться пропускна здатність~\eqref{eq:spikeentropy}. Час нарізано на клітинки довжини $\Delta t$, у кожній клітинці імпульс або є ($1$), або немає ($0$): потік імпульсів --- це двійковий рядок. Імпульс потрапляє в клітинку з імовірністю $p=r\Delta t$. Що дрібніші клітинки, то довший рядок і то більше біт; лічильник імпульсів втрачає майже все, що записано в тому, \emph{де} стоять одиниці.}
\label{fig:cells}
\end{figure}

Біти на імпульс залежать від точності часу лише логарифмічно: за точності $10\ms$ залишиться близько п'яти біт на імпульс. Але якщо отримувач лише рахує імпульси за секунду, то при $r=10$ він отримає менше двох біт за всю секунду (підрозділ~\ref{sec:rate}).

\begin{keyblock}
\begin{proposition}[Пропускна здатність імпульсного каналу]
\label{prop:capacity}
Потік імпульсів із середньою частотою $r$, прочитаний із точністю часу $\Delta t$, несе не більше~\eqref{eq:spikeentropy}. Майже вся ця пропускна здатність міститься в \emph{часі} імпульсів: отримувач, який лише рахує імпульси, видобуває з того самого потоку в десятки разів менше, ніж той, хто розрізняє їхні моменти з точністю до мілісекунди.
\end{proposition}
\end{keyblock}

Це верхня межа, а не вимірювання. Вимірювання на чутливих нейронах, для яких відомий подразник, дають від одного до кількох біт на імпульс; у найкращих випадках це до половини межі, обчисленої для їхньої точності~\cite{Rieke1997}. Отже, принаймні на вході в нервову систему час імпульсів використовується, а не пропадає.

\section{Шумний канал}

Пропускну здатність~\eqref{eq:spikeentropy} пораховано без шуму. Про те, де він виникає, твердо відомі три факти.

\emph{Сам генератор імпульсів точний.} Якщо в нейрон кори, вийнятий із мозку разом зі шматочком тканини, багато разів подати той самий струм, що швидко змінюється в часі, імпульси повторюються з точністю близько мілісекунди~\cite{Mainen1995}. Якщо струм сталий, моменти імпульсів від разу до разу розповзаються: точний час породжують швидкі перепади входу, а не сам нейрон.

\emph{Синапс ненадійний.} Імпульс, що добіг до \nt{GLUT}-синапсу, викидає порцію медіатора лише з певною ймовірністю $p$. На синапсах гіпокампа понад половина імпульсів до отримувача просто не доходить, і причина --- саме випадковість викиду, а не збої проведення~\cite{AllenStevens1994}. Для зв'язківця це канал зі стираннями; кілька синапсів між двома нейронами почасти рятують справу.

\emph{У живій корі потік імпульсів виглядає випадковим.} Інтервали між імпульсами розкидані приблизно так само, як у випадкового пуассонівського потоку, а при повторенні того самого подразника кількість імпульсів за вікно змінюється так, що її дисперсія близька до середнього~\cite{Softky1993}.

Як точний елемент породжує випадковий потік? Нейрон отримує тисячі входів, кожен ненадійний, а збудження й гальмування врівноважені, як у розділі~\ref{sec:gaba}. Напруга мембрани тремтить біля порога, і моменти імпульсів задають ці тремтіння. Шум це чи повідомлення, яке ми не вміємо прочитати, --- багато в чому відкрите питання. Частина «шуму» вже виявилася повідомленням: у зоровій корі значна частка розкиду слідує за рухами самої тварини~\cite{Stringer2019}. Складність тут принципова: добре стиснене повідомлення для того, хто не знає коду, не відрізнити від випадкового потоку.

\section{Частота: повільно, зате надійно}
\label{sec:rate}

Найпростіший код --- \emph{частотний}: число записано в тому, скільки імпульсів нейрон видає за вікно $T$. Ні стирання, ні тремтіння моментів такому кодові не страшні. Але він має ціну. Якщо потік пуассонівський, кількість імпульсів $n$ за вікно має середнє $rT$ і розкид $\sqrt{rT}$:
\begin{empheq}[box=\keybox]{equation}
\frac{\sigma_n}{\bar n} \;=\; \frac{1}{\sqrt{r\,T}} .
\label{eq:rateerror}
\end{empheq}
Щоб дізнатися частоту з точністю в десять відсотків, потрібно сто імпульсів, за двадцяти імпульсів на секунду --- п'ять секунд. Сусідні рівні розрізненні, якщо вони віддалені на пару розкидів, тому до частоти $r$ за час $T$ розрізняється близько $\sqrt{rT}$ рівнів: при $r=10$ і $T=1\,\text{с}$ це три-чотири рівні, менше двох біт.
Мозок так повільно не працює. Людина впізнає тварину на картинці, показаній на мить, і відповідь у мозку з'являється приблизно через $150\ms$~\cite{Thorpe1996}. За грубою оцінкою, за цей час сигнал проходить близько десятка послідовних ступенів обробки, по $10$--$20\ms$ на ступінь. За звичайних частот кори кожен нейрон встигає видати на ступені нуль або один імпульс, і його частота в такому вікні не визначена. Виходів два: взяти багато нейронів або дивитися на час.

\emph{Багато нейронів.} Нехай $N$ нейронів несуть те саме число, і отримувач додає їхні імпульси. Якщо їхні шуми незалежні, у~\eqref{eq:rateerror} замість $rT$ стає $NrT$: сто нейронів при $r=20$ за $15\ms$ дають тридцять імпульсів і точність близько двадцяти відсотків. Простір замінює час. Однак шуми сусідніх нейронів кори не незалежні: коефіцієнт кореляції кількостей імпульсів у пари сусідів зазвичай близько $0.1$~\cite{Zohary1994}. Якщо він дорівнює $c$, дисперсія середнього за $N$ нейронами дорівнює
\begin{empheq}[box=\keybox]{equation}
\sigma^2_N \;=\; \sigma^2_1\,\frac{1+(N-1)\,c}{N}
\;\xrightarrow[N\to\infty]{}\; c\,\sigma^2_1 .
\label{eq:correlated}
\end{empheq}

\begin{keyblock}
\begin{proposition}[Межа усереднення]
\label{prop:pooling}
За корельованого шуму усереднення за групою зменшує похибку не більше ніж у $1/\sqrt{c}$ разів. При $c=0.1$ це втричі, і такого виграшу досягають уже на кількох десятках нейронів; додавати більше марно.
\end{proposition}
\end{keyblock}

Не кожна кореляція однаково шкідлива. Обмежує точність лише та частина спільного шуму, яка \emph{схожа на сам сигнал}, тобто зсуває всю групу так само, як її зсунула б зміна вимірюваної величини~\cite{MorenoBote2014}. Скільки такого шуму в мозку насправді, ще з'ясовують.

Є й інший спосіб використати багато нейронів: записати число в тому, \emph{який} нейрон спрацював, як у визначнику напрямку на звук (рис.~\ref{fig:itd}). Це код \emph{місцем}, найнадійніший із відомих: у чутливих нейронів номер проводу каже, якої точки шкіри торкнулися чи якої висоти звук, і це встановлено багаторазово. Він дорогий за кількістю нейронів, зате швидкий: на повідомлення йде одна затримка.

\section{Час: швидко, але від чого відлічувати}

Другий вихід --- писати число в час імпульсу. Візьмімо нейрон~\eqref{eq:glutneuron} і подаймо на нього з моменту $t=0$ сталий вхід, який підняв би напругу до рівня $U$. Напруга зростає як $V=U\bigl(1-e^{-t/\tau}\bigr)$ і досягає порога $\theta$ у момент
\begin{empheq}[box=\keybox]{equation}
t_1 \;=\; \tau\,\ln\frac{U}{U-\theta}, \qquad U>\theta .
\label{eq:latency}
\end{empheq}
Сильний вхід --- ранній імпульс, слабкий --- пізній, підпороговий --- жодного. При $\tau=10\ms$ вхід $U=4\theta$ дає перший імпульс через $2.9\ms$, $U=2\theta$ --- через $6.9\ms$, $U=1.2\theta$ --- через $17.9\ms$. Один-єдиний імпульс несе число.

Але від якого моменту відлічувати $t_1$? Годинника немає, і отримувач не знає, коли почався подразник. Є три відповіді.

\emph{Відносна затримка.} Два нейрони бачать той самий подразник, і різниця їхніх затримок $t_1^A-t_1^B$ залежить лише від їхніх входів, а не від моменту початку. Порівнюють моменти детектори збігів із затримками (рис.~\ref{fig:itd}). Якщо $N$ нейронів видали по одному імпульсу, сам \emph{порядок} їхнього приходу несе до $\log_2N!$ біт: для десяти нейронів близько двадцяти двох біт, тоді як відповідь «спрацював чи ні» дала б не більше десяти. У сітківці показано, що за відносними затримками перших імпульсів її вихідних нейронів картинку можна відновити швидко й надійно~\cite{Gollisch2008}.

\emph{Залп як початок слова.} Різка зміна подразника змушує спрацювати майже одночасно безліч нейронів. Такий залп сам слугує позначкою початку, як довга пауза між словами в абетці Морзе, і отримувач відлічує затримки від нього.

\emph{Місцевий ритм.} У розділі~\ref{sec:gaba} петля \nt{GLUT}--\nt{GABA} породжувала місцевий ритм із періодом $12$--$50\ms$. Це не годинник, але всередині його циклу нейрон із сильним входом встигає спрацювати раніше, ніж нейрон зі слабким, і~\eqref{eq:latency} перетворюється на код \emph{фазою}: число записано в тому, в якій частині циклу прийшов імпульс. Такий код спостерігали: нейрони, що відповідають за певне місце в просторі, спрацьовують дедалі раніше в циклі повільного місцевого ритму в міру того, як тварина проходить через «своє» місце~\cite{OKeefe1993}. Наскільки широко мозок користується фазою, поки невідомо.

\section{Код обирає отримувач}

У послідовності імпульсів є і частота, і часи, і порядок. Що з цього буде прочитано, вирішує отримувач, і вирішує одним параметром --- своєю сталою часу $\tau$.

Усереднімо рівняння~\eqref{eq:glutneuron} за часом. Якщо на нейрон надходять незалежні потоки з частотами $r_j$, то середня напруга та її відносне тремтіння дорівнюють
\begin{empheq}[box=\keybox]{equation}
\bar V \;=\; \tau\sum_j w_j\,r_j ,
\qquad
\frac{\sigma_V}{\bar V} \;\approx\; \frac{1}{\sqrt{2\,r_{\text{вх}}\,\tau}} ,
\label{eq:reader}
\end{empheq}
де друга рівність записана для однакових ваг, а $r_{\text{вх}}$ --- сумарна частота всіх входів. Коли $\tau$ велика, напруга майже не тремтить і повторює зважену суму частот: отримувач --- \emph{інтегратор}, він читає частоти й забуває часи. Коли $\tau$ мала, напруга встигає впасти між імпульсами, і поріг досягається лише тоді, коли кілька імпульсів прийшли в межах вікна~\eqref{eq:window}: отримувач --- \emph{детектор збігів}, він читає часи (рис.~\ref{fig:reader}). Тисяча входів по п'ять імпульсів за секунду при $\tau=10\ms$ дають тремтіння близько десяти відсотків, майже чисте інтегрування. Якщо гальмування, як у розділі~\ref{sec:gaba}, вкоротило $\tau$ до $2\ms$, тремтіння зростає до двадцяти з гаком відсотків, і нейрон починає помічати збіги.

\begin{figure}[t]
\centering
\begin{tikzpicture}[>=Stealth,x=0.115cm,y=1cm]
% common input: the same spike train for both receivers
\foreach \s in {6,21,22,38,52,55.5,59,62.5,66,69.5,73,76.5,80,83.5,87,90.5,94,97.5}{\draw[thick,gray!40!black] (\s,5.25) -- (\s,5.65);}
\draw[gray!70!black] (0,5.25) -- (105,5.25);
\node[left,font=\small] at (-1,5.45) {вхід};
\node[above,font=\scriptsize,gray!40!black] at (13.5,5.65) {рідко};
\node[above,font=\scriptsize,gray!40!black] at (21.5,6.0) {пара};
\node[above,font=\scriptsize,gray!40!black] at (75,5.65) {часто};
% axes and thresholds
\foreach \y/\th in {2.7/2.5,0/1.5}{
  \draw[gray!70!black] (0,\y) -- (106,\y) node[right,black] {\small $t$};
  \draw[densely dashed,gray!60!black] (0,\y+0.6*\th) -- (104,\y+0.6*\th) node[right,black,font=\small] {$\theta$};
}
\node[anchor=south west,font=\small] at (0,4.3) {$\tau=10\ms$: інтегратор};
\node[anchor=south east,font=\small] at (104,1.0) {$\tau=2\ms$: детектор збігів};
% integrator: tau=10 ms, theta=2.5w
\begin{scope}[yshift=2.7cm]
\foreach \a/\b/\v in {0/6/0.000,6/21/1.000,21/22/1.223,22/38/2.107,38/52/1.425,52/55.5/1.351,55.5/59/1.952,59/62.5/2.376,62.5/66/0.000,66/69.5/1.000,69.5/73/1.705,73/76.5/2.201,76.5/80/0.000,80/83.5/1.000,83.5/87/1.705,87/90.5/2.201,90.5/94/0.000,94/97.5/1.000,97.5/104/1.705}{\draw[very thick,blue!60!black] plot[domain=\a:\b,samples=14] (\x,{0.6*\v*exp(-(\x-\a)/10)});}
\foreach \s/\p/\q in {6/0.000/1.000,21/0.223/1.223,22/1.107/2.107,38/0.425/1.425,52/0.351/1.351,55.5/0.952/1.952,59/1.376/2.376,62.5/1.674/2.500,62.5/2.500/0.000,66/0.000/1.000,69.5/0.705/1.705,73/1.201/2.201,76.5/1.551/2.500,76.5/2.500/0.000,80/0.000/1.000,83.5/0.705/1.705,87/1.201/2.201,90.5/1.551/2.500,90.5/2.500/0.000,94/0.000/1.000,97.5/0.705/1.705}{\draw[very thick,blue!60!black] (\s,0.6*\p) -- (\s,0.6*\q);}
\foreach \s in {62.5,76.5,90.5}{\draw[->,thick,red!70!black] (\s,1.58) -- (\s,1.95);}
\end{scope}
% coincidence: tau=2 ms, theta=1.5w
\begin{scope}[yshift=0cm]
\foreach \a/\b/\v in {0/6/0.000,6/21/1.000,21/22/1.001,22/38/0.000,38/52/1.000,52/55.5/1.001,55.5/59/1.174,59/62.5/1.204,62.5/66/1.209,66/69.5/1.210,69.5/73/1.210,73/76.5/1.210,76.5/80/1.210,80/83.5/1.210,83.5/87/1.210,87/90.5/1.210,90.5/94/1.210,94/97.5/1.210,97.5/104/1.210}{\draw[very thick,blue!60!black] plot[domain=\a:\b,samples=14] (\x,{0.6*\v*exp(-(\x-\a)/2)});}
\foreach \s/\p/\q in {6/0.000/1.000,21/0.001/1.001,22/0.607/1.500,22/1.500/0.000,38/0.000/1.000,52/0.001/1.001,55.5/0.174/1.174,59/0.204/1.204,62.5/0.209/1.209,66/0.210/1.210,69.5/0.210/1.210,73/0.210/1.210,76.5/0.210/1.210,80/0.210/1.210,83.5/0.210/1.210,87/0.210/1.210,90.5/0.210/1.210,94/0.210/1.210,97.5/0.210/1.210}{\draw[very thick,blue!60!black] (\s,0.6*\p) -- (\s,0.6*\q);}
\foreach \s in {22}{\draw[->,thick,red!70!black] (\s,0.98) -- (\s,1.35);}
\end{scope}
\foreach \x in {0,50,100}{\draw[gray!60!black] (\x,-0.06) -- (\x,0.06) node[below,font=\scriptsize] {\x\,мс};}
\end{tikzpicture}
\caption{Той самий вхід --- два різні отримувачі. Кожен імпульс піднімає напругу на $w$, після чого вона стікає, як у моделі нейрона~\eqref{eq:glutneuron}; червоні стрілки --- імпульси отримувача. Повільний отримувач ($\tau=10\ms$, $\theta=2.5w$) спрацьовує там, де імпульсів \emph{багато}, і не помічає пари. Швидкий ($\tau=2\ms$, $\theta=1.5w$) спрацьовує лише на пару в межах вікна~\eqref{eq:window}, а частий, але не збіжний потік пропускає.}
\label{fig:reader}
\end{figure}

Вимірювання показують, що в активній корі провідність мембрани зростає в кілька разів порівняно зі спокоєм~\cite{Destexhe2003}. Отже, $\tau$ там коротша, і нейрони кори в робочому режимі ближчі до детекторів збігів, ніж до інтеграторів.

\begin{keyblock}
\begin{proposition}[Код задає отримувач]
\label{prop:reader}
Та сама послідовність імпульсів для повільного отримувача несе частоту, для швидкого --- часи, для об'єму, у який викинуто медіатор, --- рівень концентрації, згладжений за секунди (розділ~\ref{sec:serocomputer}). Який код використовується, визначає не відправник, а стала часу отримувача, і її підлаштовує гальмування.
\end{proposition}
\end{keyblock}
\section{Крапки й тире: синапс як фільтр}

Досі синапс був у нас проводом із вагою. Насправді його відповідь залежить від попередніх імпульсів, і це дає мові нейронів другий знак.

\emph{Виснаження: синапс передає зміни.} Кожен викид витрачає частку $U$ запасу медіатора, готового до викиду~\cite{Tsodyks1997}, а запас відновлюється зі сталою часу $\tau_{\text{відн}}$ порядку сотень мілісекунд. За частоти $r$ амплітуда відповіді на імпульс установлюється приблизно на рівні $A(r)=A_0/(1+U r\tau_{\text{відн}})$, де $A_0$ --- відповідь відпочилого синапсу. Струм, який отримувач отримує за одиницю часу, дорівнює
\begin{empheq}[box=\keybox]{equation}
r\,A(r) \;=\; \frac{A_0\,r}{1+U r\,\tau_{\text{відн}}}
\;\xrightarrow[r\to\infty]{}\; \frac{A_0}{U\tau_{\text{відн}}} .
\label{eq:depression}
\end{empheq}
При $U=0.5$ і $\tau_{\text{відн}}=0.8\,\text{с}$ насичення починається вже з $1/(U\tau_{\text{відн}})=2.5$ імпульсу за секунду. Якщо частота зросла з п'яти до двадцяти, усталений струм зросте лише на третину. Зате перші імпульси на новій частоті приходять, поки запас ще колишній, і струм на мить підскакує вчетверо, а потім за $\tau_{\text{відн}}/(1+Ur\tau_{\text{відн}})\approx90\ms$ спадає (рис.~\ref{fig:depression}).

Такий синапс передає не частоту, а її \emph{зміну}, по суті --- відносну зміну $\Delta r/r$~\cite{Abbott1997depression}. Для електронника це фільтр верхніх частот, поставлений просто в провід.

\begin{figure}[t]
\centering
\begin{tikzpicture}[>=Stealth,x=12.5cm,y=1cm]
% input rate
\draw[->,gray!70!black] (0,2.6) -- (0.9,2.6) node[right,black] {\small $t$};
\draw[->,gray!70!black] (0,2.6) -- (0,4.3) node[above,black] {\small $r$};
\draw[very thick,blue!60!black] (0,2.9) -- (0.2,2.9) -- (0.2,3.8) -- (0.8,3.8);
\node[left,font=\scriptsize] at (0,2.9) {$5$};
\node[left,font=\scriptsize] at (0,3.8) {$20$};
\node[right,font=\small,blue!60!black] at (0.4,4.1) {частота на вході синапсу};
% output current
\draw[->,gray!70!black] (0,0) -- (0.9,0) node[right,black] {\small $t$};
\draw[->,gray!70!black] (0,0) -- (0,2.45) node[left,black] {\small $rA$};
\draw[densely dashed,gray!60!black] (0,0.667) -- (0.8,0.667);
\draw[very thick,red!70!black] (0,0.5) -- (0.2,0.5) -- (0.2,2.0)
   plot[domain=0.2:0.8,samples=60] (\x,{0.667+(2.0-0.667)*exp(-(\x-0.2)/0.089)});
\node[left,font=\scriptsize] at (0,0.5) {$1.7$};
\node[left,font=\scriptsize] at (0,2.0) {$6.7$};
\node[right,font=\scriptsize] at (0.8,0.667) {$2.2$};
\node[right,font=\small,red!70!black] at (0.28,1.6) {струм до отримувача: підскок і спад за $\approx90\ms$};
\draw[gray!60!black] (0.2,-0.08) -- (0.2,0.08); \node[below,font=\scriptsize] at (0.2,0) {$0.2\,\text{с}$};
\draw[gray!60!black] (0.8,-0.08) -- (0.8,0.08); \node[below,font=\scriptsize] at (0.8,0) {$0.8\,\text{с}$};
\end{tikzpicture}
\caption{Синапс, що виснажується, за формулою~\eqref{eq:depression} при $U=0.5$, $\tau_{\text{відн}}=0.8\,\text{с}$; струм в одиницях $A_0$ за секунду, штрихова лінія --- усталений струм: він зріс лише з $1.7$ до $2.2$, а в момент стрибка --- до $6.7$. Синапс повідомляє отримувачеві про те, що частота змінилася, а не про те, яка вона.}
\label{fig:depression}
\end{figure}

\emph{Полегшення: пачка як тире.} В інших синапсів імовірність викиду мала, але після кожного імпульсу на десятки мілісекунд зростає. Одиночний імпульс через такий синапс найчастіше губиться. А \emph{пачка} --- кілька імпульсів поспіль з інтервалами в кілька мілісекунд --- проходить майже напевно. Навіть без полегшення ймовірність, що пропадуть усі $k$ імпульсів пачки, дорівнює
\begin{empheq}[box=\keybox]{equation}
P_{\text{втрати}} \;=\; (1-p)^k ,
\label{eq:burst}
\end{empheq}
при $p=0.2$ це $0.8$ для одиночного імпульсу й $0.41$ для пачки з чотирьох; полегшення зменшує її ще більше. Так у нейрона з'являється другий знак: одиночний імпульс --- крапка, пачка --- тире. Синапс, що виснажується, найкраще передає перший імпульс після паузи; полегшуваний пропускає пачки й глушить одиночні крапки. Те, що пачки передаються надійніше, виміряно; що мозок справді користується ними як окремим знаком, --- правдоподібна гіпотеза~\cite{Lisman1997}.

\section{Як говорять \nt{GABA}-нейрони}

Найчисленніші з них у корі --- швидкі~\cite{Tremblay2016}. Їхні імпульси коротші, ніж у \nt{GLUT}-нейронів, вони можуть видавати їх рівно й довго з частотою в сотні за секунду, майже не втомлюючись. Їхні синапси надійніші: зв'язок із кожним отримувачем складається з багатьох місць викиду, й імпульс майже ніколи не губиться. Їхні виходи сідають поруч із тим місцем отримувача, де народжується його імпульс, тож вони керують не тим, як отримувач додає входи, а тим, чи спрацює він і коли.

Самі вони отримують збудження від багатьох сусідніх \nt{GLUT}-нейронів і повідомляють про сумарну активність навколо --- саме те, що потрібно для ділення з розділу~\ref{sec:gaba}. Нарешті, швидкі \nt{GABA}-нейрони зв'язані один з одним і легко спрацьовують разом; саме вони задають місцевий ритм, про який ішлося вище~\cite{Cardin2009}.

Мовою абетки Морзе \nt{GABA}-нейрони відповідають не за літери, а за \emph{паузи}. Вони ріжуть неперервний потік \nt{GLUT}-імпульсів на вікна, усередині яких отримувач може спрацювати, звужують вікна збігу й приглушують увесь потік, коли він стає надто гучним.

\begin{keyblock}
\begin{proposition}[Поділ ролей]
\label{prop:punctuation}
\nt{GLUT}-нейрони несуть зміст повідомлення: \emph{що} і \emph{скільки}. Швидкі \nt{GABA}-нейрони несуть його розмітку: \emph{коли} можна говорити й \emph{наскільки гучно}; ще один клас, гальмуючи гальмівних, вирішує, \emph{для кого} ворота відчинені. Окремі частини цього поділу виміряно; як загальне правило це робоча гіпотеза.
\end{proposition}
\end{keyblock}

Є й інші, повільніші \nt{GABA}-нейрони, виходи яких сідають на окремі входи отримувача. Вони працюють як заборона~\eqref{eq:veto} з розділу~\ref{sec:gaba}: закривають один потік \nt{GLUT}-імпульсів, не чіпаючи решти.

Поділ на швидкі й повільні \nt{GABA}-нейрони --- стара картина, і вона неповна. Зараз у корі розрізняють щонайменше три великі класи~\cite{Tremblay2016}, і третій влаштований несподівано: він гальмує не \nt{GLUT}-нейрони, а інші \nt{GABA}-нейрони, насамперед ті, що закривають входи~\cite{Pfeffer2013}. Гальмування гальмування --- подвійне заперечення. Такий нейрон \emph{відчиняє} ворота, не посилаючи отримувачеві жодного збуджувального імпульсу. Вмикають його сигнали з інших областей кори та широкомовні канали, тож він працює як дозвільний вхід цілої ділянки мережі: поки він активний, заборони знято, і потік проходить~\cite{Pi2013}. У додатку~\ref{sec:othermediators} цей прийом названо розгальмовуванням.
\section{Ощадлива мова}

Останнє, що відомо про мову нейронів твердо, --- вона дорога. Імпульс разом із роботою синапсів, яку він викликає, обходиться приблизно в $10^9$ молекул АТФ (оцінка для кори гризунів~\cite{Attwell2001}), а одна молекула дає близько $10^{-19}$ джоуля. Отже, імпульс коштує порядку $E_{\text{імп}}\sim10^{-10}$ джоуля. Увесь мозок споживає близько $20$ ватів, і якщо на передачу сигналів іде половина, $P\sim10$ ватів, то середня частота одного нейрона
\begin{empheq}[box=\keybox]{equation}
\bar r \;\approx\; \frac{P}{N\,E_{\text{імп}}}
\;\sim\; \frac{10\ \text{Вт}}{8.6\cdot10^{10}\cdot10^{-10}\ \text{Дж}}
\;\approx\; 1\ \text{імпульс за секунду}.
\label{eq:energy}
\end{empheq}
Докладніші оцінки для кори дають ще менше~\cite{Lennie2003}. Код мозку мусить бути \emph{розрідженим}: швидко працювати може лише мала частка нейронів.

З~\eqref{eq:spikeentropy} видно, що це не так уже й погано. За точності $1\ms$ імпульс нейрона, що працює з частотою $100$ за секунду, несе не більше п'яти біт, а імпульс нейрона, що працює раз на секунду, --- до одинадцяти. Якщо платити доводиться за імпульси, вигідно говорити рідко. Кора й справді міняє точність на енергію: за нестачі їжі нейрони зберігають частоту імпульсів, але витрачають менше на синаптичні струми, і їхні відповіді стають менш точними~\cite{Padamsey2022}.

В абетці Морзе часті літери короткі: найчастіша літера англійського тексту E --- одна крапка. Код нейронів, наскільки відомо, дотримується того самого правила в іншій формі: те, що очікуване, коштує мало імпульсів~\cite{Barlow1961}. Нейрон за сталого подразника поступово знижує частоту, датчики розділу~\ref{sec:glutcomputer} відповідають на зміни, а не на рівні, а \nt{DOPA}-нейрони розділу~\ref{sec:neurontypes} повідомляють лише про помилку передбачення винагороди.

\begin{keyblock}
\begin{proposition}[Ощадливість]
\label{prop:economy}
Енергія обмежує середню частоту нейрона величиною порядку одного імпульсу за секунду. Тому мова \nt{GLUT}-нейронів розріджена й витрачає імпульси на новини: на те, що змінилося або не було передбачене. Енергетичну межу виміряно; що код мозку підлаштований під найбільшу кількість біт на джоуль, --- добре обґрунтована гіпотеза.
\end{proposition}
\end{keyblock}

\section{Підсумок}

\begin{table}[t]
\centering
\footnotesize
\setlength{\tabcolsep}{4pt}
\renewcommand{\arraystretch}{1.15}
\begin{tabular}{>{\raggedright}p{2.6cm}>{\raggedright}p{2.3cm}>{\raggedright}p{3.0cm}>{\raggedright}p{2.0cm}>{\raggedright\arraybackslash}p{3.6cm}}
\toprule
Спосіб запису & Що несе & Хто читає & Час на повідомлення & Що відомо \\
\midrule
номер нейрона, що спрацював (код місцем) & яке значення & будь-який отримувач; детектори збігів & одна затримка & встановлено в чутливих нейронів і у визначенні напрямку на звук \\
частота одного нейрона & величину & інтегратор із великою $\tau$; об'єм (розд.~\ref{sec:serocomputer}) & сотні мс -- секунди & встановлено; для ступеня обробки в $10$--$20\ms$ надто повільно \\
частота групи & величину & нейрон із тисячами входів & $10$--$50\ms$ & встановлено; виграш обмежений кореляціями~\eqref{eq:correlated} \\
час першого імпульсу, порядок & величину, порядок & детектори збігів із затримками & одна затримка & встановлено на вході нервової системи; у корі обговорюється \\
фаза в місцевому ритмі & величину, положення & нейрони зі спільним ритмом & цикл $12$--$50\ms$ і повільніше & спостерігається в окремих областях; загальна роль --- гіпотеза \\
пачка («тире») & «важливо»: надійна доставка & полегшуваний синапс & $\sim10\ms$ & надійність виміряно; окремий знак --- гіпотеза \\
зміна частоти & новину & синапс, що виснажується & $\sim0.1\,\text{с}$ & встановлено \\
імпульси швидких \nt{GABA}-нейронів & коли й наскільки гучно & усі сусідні \nt{GLUT}-нейрони & $1$--$30\ms$ & ритм і ділення виміряно; «пунктуація» як правило --- гіпотеза \\
\bottomrule
\end{tabular}
\caption{Способи, якими \nt{GLUT}- і \nt{GABA}-нейрони записують повідомлення імпульсами. Правий стовпець відділяє виміряне від припущеного.}
\label{tab:code}
\end{table}

Таблиця~\ref{tab:code} збирає те, що ми знаємо; видно з неї й те, чого ми не знаємо. Ми не знаємо, наскільки кора користується точним часом імпульсів, а не лише частотами груп. Ми не знаємо, чи є в нейронів «слова» --- стійкі поєднання імпульсів багатьох нейронів, які читаються як ціле. І ми не знаємо, чи однакова мова в різних частинах мозку. Абетку Морзе можна вивчити за вечір, бо її таблиця відома. Таблицю мови нейронів ще належить скласти, і складатимуть її, найімовірніше, інженери зв'язку.

\section{Задачі}

\begin{exercise}[\lvlA]
\label{ex:04:1}
\emph{Пропускна здатність у числах.} $\Delta t=1\ms$. (а)~Скільки біт на джоуль дає нейрон із частотою $1$ імпульс за секунду за ціни імпульсу з~\eqref{eq:energy}? (б)~При $r=10$ у скільки разів менше за межу~\eqref{eq:spikeentropy} видобуває отримувач, який лише рахує імпульси за секунду?
\end{exercise}

\begin{exercise}[\lvlB]
\label{ex:04:2}
\emph{Оптимальна частота.} Нейрон витрачає потужність $P_0+rE_{\text{імп}}$, де $P_0$ --- друга половина з $20$~Вт, поділена на всі нейрони, $E_{\text{імп}}=10^{-10}$~Дж. За якої частоти $r$ кількість біт на джоуль $R_{\max}(r)/(P_0+rE_{\text{імп}})$ за~\eqref{eq:spikeentropy} з $\Delta t=1\ms$ найбільша?
\end{exercise}

\begin{exercise}[\lvlB]
\label{ex:04:3}
\emph{Яка кореляція шкідлива.} $N=1000$ нейронів, дисперсія $\sigma^2=1$, попарна кореляція $c=0.1$. Обчисліть інформацію Фішера $\mathcal I=f'^{\mathsf T}\Sigma^{-1}f'$ ($\Sigma$ --- матриця коваріацій) для однакових нахилів $f'=\mathbf 1$ і для нахилів $\pm1$ з $\mathbf 1^{\mathsf T}f'=0$. У скільки разів вони різняться?
\end{exercise}

\begin{exercise}[\lvlB]
\label{ex:04:4}
\emph{Моделювання: детектор збігів.} Нейрон~\eqref{eq:glutneuron} отримує $1000$ пуассонівських входів по $5$ імпульсів за секунду, $w=1$; поріг $\theta$ такий, що вільна напруга перетинає його раз на секунду. Моделлю знайдіть, скільки одночасних входів $m$ запалюють нейрон з імовірністю $1/2$ при $\tau=10$ і $2\ms$.
\end{exercise}

\begin{exercise}[\lvlB]
\label{ex:04:5}
\emph{Проєктування: детектор відносних змін.} Доберіть синапс~\eqref{eq:depression}: усталений струм при $40$ імпульсах за секунду не більш ніж на $10\,\%$ вищий, ніж при $10$, а після стрибка до $40$ струм виходить на новий рівень зі сталою часу не більше $50\ms$. Яке найменше $\tau_{\text{відн}}$ і за якого $U$?
\end{exercise}

\begin{exercise}[\lvlA]
\label{ex:04:6}
\emph{Час першого імпульсу й пачки.} (а)~За~\eqref{eq:latency} знайдіть вхід, за якого перший імпульс приходить рівно через одну сталу часу. Від чого він залежить? (б)~За ймовірності викиду $p=0.2$ скільки імпульсів потрібно в пачці, щоб утратити всі з імовірністю нижче $5\,\%$?
\end{exercise}

\begin{exercise}[\lvlC]
\label{ex:04:7}
\emph{Дослідження: скільки біт у розташуванні.} Нейрон --- незалежні клітинки~\eqref{eq:spikeentropy}, $r=10$, $\Delta t=1\ms$. (а)~Розкладіть ентропію «слова» з $20$ клітинок на ентропію кількості імпульсів і залишок. (б)~Змоделюйте оцінку ентропії за частотами слів із $N=30$ і $10^4$ записів. Наскільки вона занижена?
\end{exercise}
